\section{Common learning in arbitrary input dimension}
\label{sec:matrixlearning77}

The two-dimensional learner extends to the full effect body through a
calibration step.  The calibration must resolve the two support
endpoints simultaneously.  A bound in operator norm alone does not
provide this information: after compression, a poorly aligned
subspace can introduce noise at a scale larger than $N^{-1}$.  The
procedure below learns a basis in which this additional noise is
bounded by the finite-horizon cutoff.

Throughout this section, $\mathcal M_E$ is the ordered binary
measurement associated with $E\in\mathfrak E_d$.  Training may use
entangled input blocks and classical feedback, but it has no access
to the measurement environment.  Its output is a classical
description of a legal effect.  The future horizon $N$ and the number
of training calls remain separate resources.

\begin{theorem}[Common learning on the matrix effect body]
\label{thm:matrixlearning77}
There are absolute constants $c,C,\delta_0>0$ with the following
property.  For every $d,N\ge1$, $0<\delta\le\delta_0$, and
$0<\eta\le1/8$, one common procedure returns a legal effect
$\widehat E\in\mathfrak E_d$ such that
\begin{equation}\label{eq:matrixlearningrisk77}
 \sup_{E\in\mathfrak E_d}
 \mathbb P_E\{d_N(\mathcal M_E,\mathcal M_{\widehat E})>\delta\}
 \le\eta.
\end{equation}
On every record its number of calls satisfies
\begin{equation}\label{eq:matrixlearningcalls77}
 M\le C d^4N\delta^{-2}\log(d/\eta),
\end{equation}
and every entangled input block has length at most $N$.  Conversely,
any common procedure satisfying \eqref{eq:matrixlearningrisk77}
with at most $M$ calls on every record satisfies
\begin{equation}\label{eq:matrixlearninglower77}
 M\ge cN\delta^{-2}\log(1/\eta).
\end{equation}
The converse permits arbitrary reference-assisted adaptive training
and bounded public stopping.  In particular, for each fixed $d$, the
minimax number of calls is
$\Theta_d(N\delta^{-2}\log(1/\eta))$ on the complete closed effect
body.  The displayed dependence on $d$ in the upper bound is not
asserted to be optimal.
\end{theorem}

We first give the calibration, including its finite budgets.  All
matrix orders below are Loewner orders.  The use of a basis selected
from earlier observations is always followed by fresh samples.

\begin{lemma}[Simultaneous endpoint calibration]
\label{lem:matrixcalibration77}
Let $\tau=N^{-1}$ and $0<\epsilon\le1/8$.  There is a common
product-input procedure returning $A\in\mathfrak E_d$ with at most
$Cd^4N\log(d/\epsilon)$ calls on every record.  With probability at
least $1-\epsilon$, it satisfies
\begin{align}
 \tfrac12(A+\tau I)&\preceq E+\tau I\preceq2(A+\tau I),
                                      \label{eq:matrixcalibrationplus77}\\
 \tfrac12(I-A+\tau I)&\preceq I-E+\tau I
                                  \preceq2(I-A+\tau I),
                                      \label{eq:matrixcalibrationminus77}\\
 \|E-A\|_{\rm op}&\le\sqrt\tau.
                                      \label{eq:matrixcalibrationnorm77}
\end{align}
The input states depend only on the public parameters and the
preceding record; no spectral data about $E$ are supplied.
\end{lemma}
\begin{proof}
Set $J=\lceil\log_2N\rceil$, $H=2^J$, and $t_j=2^{-j}$ for
$0\le j\le J$.  Thus $N\le H<2N$ unless $N=1$, and
$\tau/2<t_J\le\tau$.  Start from $A_0=I/2$.  If $J=0$, return
$A_0$: all three conclusions hold without samples.

We use
\[
 h_d=d+4\binom d2=2d^2-d,
 \qquad K_d=(320d)^2,
 \qquad \alpha=1/16.
\]
At stage $j\ge1$, assign failure allowance and sample size
\begin{equation}\label{eq:matrixcalibrationbudget77}
 \epsilon_j=\frac{\epsilon 2^j}{4H},\qquad
 L_j=\log\frac{2h_d}{\epsilon_j},\qquad
 n_j=\left\lceil K_d t_j^{-1}L_j\right\rceil.
\end{equation}
Choose an orthonormal eigenbasis $e_1,\ldots,e_d$ of the current
legal estimate $A_{j-1}$.  The choice is deterministic and measurable,
including at multiple eigenvalues.  More explicitly, fix a public
input basis.  All calibration matrices have algebraic entries in that
basis: the sample frequencies are rational, and spectral projections,
clipping at rational endpoints, and normalization preserve algebraic
entries.  Order the distinct eigenvalues, project the public basis
onto each eigenspace, and apply Gram--Schmidt in that fixed order,
discarding exactly zero vectors.  Exact algebraic arithmetic
therefore supplies a finite choice at every stage and requires no
oracle for arbitrary real spectral data.
Make $n_j$ independent calls on each of the $h_d$ states
\begin{equation}\label{eq:matrixcalibrationstates77}
 e_i,\qquad
 \frac{e_i+e_j}{\sqrt2},\quad
 \frac{e_i-e_j}{\sqrt2},\quad
 \frac{e_i+\mathrm i e_j}{\sqrt2},\quad
 \frac{e_i-\mathrm i e_j}{\sqrt2}
 \quad(i<j).
\end{equation}
Here and below a unit vector denotes its pure input state.  Estimate
each diagonal entry of $E$ by its empirical outcome frequency.
Estimate the real part of each off-diagonal entry by half the
frequency difference between the positive and negative real
superpositions.  Its imaginary part is half the frequency on
$(e_i-\mathrm i e_j)/\sqrt2$ minus half the frequency on
$(e_i+\mathrm i e_j)/\sqrt2$.  These formulas define a Hermitian
matrix $B_j$.

To define the procedure on every record, clip the eigenvalues of
$B_j$ to $[-\alpha t_j,1+\alpha t_j]$, obtaining $\widetilde B_j$,
and set
\begin{equation}\label{eq:matrixcalibrationupdate77}
 A_j=\frac{\widetilde B_j+\alpha t_jI}{1+2\alpha t_j}.
\end{equation}
The returned matrix is always a legal effect.  In particular, a
failed concentration event cannot affect the subsequent call budget
or require an additional experiment.

We prove the invariant
\begin{equation}\label{eq:matrixcalibrationinvariant77}
 \tfrac12(A_j+t_jI)\preceq E+t_jI\preceq2(A_j+t_jI),
 \qquad
 \tfrac12(I-A_j+t_jI)\preceq I-E+t_jI
                                      \preceq2(I-A_j+t_jI)
\end{equation}
on the successive good events.  It holds for $j=0$ because
$0\preceq E\preceq I$.

Fix a preceding good record and abbreviate
$t=t_j$, $A=A_{j-1}$, $B=B_j$, $D=B-E$, and $q=L_j/n_j$.
All expectations and probabilities in the next calculation are
conditional on that record.  Bernstein's inequality and a union
bound over the $h_d$ groups show, with failure at most
$2h_de^{-L_j}=\epsilon_j$, that every empirical frequency has error
at most
\[
 \sqrt{2p(1-p)q}+2q/3.
\]
For a pair of opposite superposition states, the sum of the outcome
probabilities is $E_{ii}+E_{jj}$.  Cauchy--Schwarz in each
half-difference therefore gives the following two simultaneous
bounds, also valid for $i=j$ after enlarging the constants:
\begin{equation}\label{eq:matrixcalibrationentry77}
 |D_{ij}|\le\sqrt{2(E_{ii}+E_{jj})q}+2q,
 \qquad
 |D_{ij}|\le\sqrt{2(2-E_{ii}-E_{jj})q}+2q.
\end{equation}
For an off-diagonal complex entry, apply the real estimate to the
two quadratures and take their Euclidean norm.  The factor $\sqrt2$
in \eqref{eq:matrixcalibrationentry77} includes this operation.

Write $a_i$ for the eigenvalues of $A$.  The preceding invariant,
whose regularization is $2t$, gives
\[
 E_{ii}\le2a_i+2t,
 \qquad 1-E_{ii}\le2(1-a_i)+2t,
 \qquad q\le t/K_d.
\]
Since
\[
 t(a_i+a_j+2t)\le2(a_i+t)(a_j+t),
 \qquad t\le\sqrt{(a_i+t)(a_j+t)},
\]
the first bound in \eqref{eq:matrixcalibrationentry77} yields
\begin{equation}\label{eq:matrixcalibrationnormalized77}
 \frac{|D_{ij}|}{\sqrt{(a_i+t)(a_j+t)}}
 \le\frac{2\sqrt2}{\sqrt{K_d}}+\frac2{K_d}
 \le\frac5{\sqrt{K_d}}.
\end{equation}
Replacing $a_i$ by $1-a_i$ gives the same estimate for the
complement.  The Hilbert--Schmidt norm bounds the operator norm, so
our choice of $K_d$ implies
\begin{align}
 -\tfrac1{64}(A+tI)&\preceq D\preceq\tfrac1{64}(A+tI),\notag\\
 -\tfrac1{64}(I-A+tI)&\preceq D
                            \preceq\tfrac1{64}(I-A+tI).
                                      \label{eq:matrixcalibrationraworder77}
\end{align}
The preceding invariant and $t_{j-1}=2t$ also imply
\[
 A+tI\preceq4(E+tI),\qquad
 I-A+tI\preceq4(I-E+tI).
\]
Consequently
\begin{equation}\label{eq:matrixcalibrationtrueorder77}
 -\alpha(E+tI)\preceq D\preceq\alpha(E+tI),
 \qquad
 -\alpha(I-E+tI)\preceq D\preceq\alpha(I-E+tI).
\end{equation}
These inequalities place the spectrum of $B$ inside
$[-\alpha t,1+\alpha t]$.  Thus the clipping in
\eqref{eq:matrixcalibrationupdate77} does nothing on this event.
On it,
\[
 A_j-E=\frac{D+\alpha t(I-2E)}{1+2\alpha t}.
\]
Since $-I\preceq I-2E\preceq I$, equations
\eqref{eq:matrixcalibrationtrueorder77} imply
\begin{align*}
 -2\alpha(E+tI)&\preceq A_j-E\preceq2\alpha(E+tI),\\
 -2\alpha(I-E+tI)&\preceq A_j-E
                                      \preceq2\alpha(I-E+tI).
\end{align*}
Adding the corresponding regularized endpoint gives factors
$1-2\alpha=7/8$ and $1+2\alpha=9/8$.  This proves the weaker
factor-two invariant \eqref{eq:matrixcalibrationinvariant77} for
stage $j$.

The unweighted entry estimate also gives the operator-norm bound
needed later.  Because $E_{ii}+E_{jj}\le2$, we have
\[
 \|D\|_{\rm op}\le4d\sqrt{t/K_d}=\sqrt t/80.
\]
It follows that
\begin{equation}\label{eq:matrixcalibrationstageerror77}
 \|A_j-E\|_{\rm op}
 \le\sqrt t/80+\alpha t\le\sqrt t.
\end{equation}
At the final stage, adding $(\tau-t_J)I$ preserves the factor-two
orders, and \eqref{eq:matrixcalibrationstageerror77} gives
\eqref{eq:matrixcalibrationnorm77}.

The conditional failure allowances sum to less than $\epsilon/2$,
since $\sum_{j=1}^J2^j<2H$.  This proves the asserted probability
bound.  Finally, the deterministic call count is $h_d\sum_j n_j$.
Writing $2^j=H/2^k$ yields
\[
 \sum_{j=1}^J2^jL_j
 \le H\sum_{k=0}^{J-1}2^{-k}
       \left(\log\frac{8h_d}{\epsilon}+k\log2\right)
 \le CH\log(d/\epsilon).
\]
The ceiling terms contribute at most $h_dJ\le h_dH$.
As $h_dK_d\le Cd^4$ and $H<2N$ when $N>1$, the stated bound
holds on every record.  In particular, the number of dyadic stages
does not introduce a factor $\log N$.
\end{proof}

\begin{lemma}[Calibrated variance and compression]
\label{lem:matrixcompression77}
Suppose $A,E\in\mathfrak E_d$ satisfy
\eqref{eq:matrixcalibrationplus77}--\eqref{eq:matrixcalibrationnorm77}.
Set $W_G=G(I-G)+\tau I/2$.  Then
\begin{equation}\label{eq:matrixcalibratedvariance77}
 \tfrac1{12}W_A\preceq W_E\preceq12W_A.
\end{equation}
If $P$ is any orthogonal projection commuting with $A$ and
$C=PEP|_{P\C^d}$, then, on $P\C^d$,
\begin{equation}\label{eq:matrixcompressionvariance77}
 \tfrac1{12}PW_AP\preceq W_C\preceq14PW_AP.
\end{equation}
All constants are independent of the spectrum, its multiplicities,
and the rank of $P$.
\end{lemma}
\begin{proof}
For a legal effect $G$, define the harmonic mean without a factor
of two by
\[
 R_G=\bigl((G+\tau I)^{-1}+(I-G+\tau I)^{-1}\bigr)^{-1}.
\]
The two factors commute because both are functions of $G$, so
\begin{equation}\label{eq:matrixharmonicidentity77}
 R_G=\frac{G(I-G)+(\tau+\tau^2)I}{1+2\tau},
 \qquad
 \tfrac13\bigl(G(I-G)+\tau I\bigr)
       \preceq R_G\preceq G(I-G)+\tau I.
\end{equation}
The inequalities use only $0<\tau\le1$.
Inversion reverses Loewner order.  Apply it to both endpoint
comparisons, add, and invert again to get
$R_E\asymp R_A$ with factors two.  Equation
\eqref{eq:matrixharmonicidentity77} therefore compares
$E(I-E)+\tau I$ and $A(I-A)+\tau I$ with factors six.
Replacing the regularization $\tau I$ by $\tau I/2$ proves
\eqref{eq:matrixcalibratedvariance77}.

The compression identity is exact:
\begin{equation}\label{eq:matrixcompressionidentity77}
 C(I-C)=PE(I-E)P+PE(I-P)EP.
\end{equation}
Because $P$ commutes with $A$,
$PE(I-P)=P(E-A)(I-P)$.  Thus
\[
 0\preceq PE(I-P)EP
       \preceq\|E-A\|_{\rm op}^2P\preceq\tau P.
\]
The first term in \eqref{eq:matrixcompressionidentity77}, after
adding $\tau P/2$, is $PW_EP$.  Apply
\eqref{eq:matrixcalibratedvariance77} and use
$PW_AP\succeq\tau P/2$ to absorb the final $\tau P$ with a
factor two.  This gives \eqref{eq:matrixcompressionvariance77}.
\end{proof}

\begin{proof}[Proof of Theorem~\ref{thm:matrixlearning77}]
For $d=1$, Bernoulli sampling and
Lemma~\ref{lem:learnendpoint76} give the upper bound directly.
Assume $d\ge2$ and put $s=\binom d2$.  Apply
Lemma~\ref{lem:matrixcalibration77} with failure allowance $\eta/2$.
Conditional on its good event, fix an orthonormal eigenbasis of
the returned $A$, with eigenvalues $a_i$, and write
\[
 w_i=a_i(1-a_i)+\tau/2.
\]
For each $i<j$, let $P_{ij}$ project onto
$\operatorname{span}(e_i,e_j)$ and let $C_{ij}$ be the corresponding
compression of $E$.  An input qubit, including its entanglement
with other input qubits or a retained reference, can be embedded
isometrically into this two-dimensional subspace.  The resulting
observed channel is exactly $\mathcal M_{C_{ij}}$.

Run the common qubit learner of
Theorem~\ref{thm:commonlearn76} on each compression, using fresh
records, accuracy
\begin{equation}\label{eq:matrixlearningpairaccuracy77}
 \varepsilon=c_*\delta/d,
\end{equation}
and failure allowance $\eta/(2s)$.  The constant $c_*>0$ will be
chosen below.  This returns legal $2\times2$ effects $\widehat C_{ij}$.
Every guarantee is conditional on the calibration record and is
uniform over the resulting compression.  A union bound therefore
shows that, with total failure at most $\eta$, calibration succeeds
and
\begin{equation}\label{eq:matrixlearningpairrisk77}
 d_N(\mathcal M_{C_{ij}},\mathcal M_{\widehat C_{ij}})
 \le\varepsilon\quad\hbox{for every }i<j.
\end{equation}

We record explicitly how an operational error controls an endpoint
form.  Theorem~\ref{thm:matrixmetric76} in dimension two and
Lemma~\ref{lem:matrixlocal76} imply that there are absolute
$a_0,A_0>0$ such that, for legal $C,D\in\mathfrak E_2$,
\begin{equation}\label{eq:matrixlearninglocal77}
 d_N(\mathcal M_C,\mathcal M_D)\le u\le a_0
 \quad\Longrightarrow\quad g_{N,C}(D-C)\le A_0u.
\end{equation}
Indeed, for $u<1/16384$ the lower metric bound forces
$Q_N(C,D)\le16384u$; decreasing $a_0$ makes this at most $1/4$,
where the local comparison applies.  Choose $c_*$ small enough that
$\varepsilon\le a_0$ and lies within the qubit learner's accuracy
range.

For a Hermitian matrix $X$ on $\operatorname{span}(e_i,e_j)$ put
\[
 \gamma_{ij}(X)^2
 =N\left(\frac{|X_{ii}|^2}{2w_i}
          +\frac{|X_{jj}|^2}{2w_j}
          +\frac{2|X_{ij}|^2}{w_i+w_j}\right).
\]
Equation \eqref{eq:matrixcompressionvariance77}, applied to left
plus right multiplication and then inverted, gives
\[
 \gamma_{ij}(X)\le\sqrt{14}\,g_{N,C_{ij}}(X).
\]
Hence \eqref{eq:matrixlearningpairrisk77} and
\eqref{eq:matrixlearninglocal77} imply
\begin{equation}\label{eq:matrixlearningpairform77}
 \gamma_{ij}(\widehat C_{ij}-C_{ij})
 \le\sqrt{14}A_0\varepsilon.
\end{equation}

Assemble a Hermitian matrix $B$ by taking each off-diagonal entry
from its corresponding estimate $\widehat C_{ij}$.  For each
diagonal entry choose once and for all one incident pair and use
that pair's estimate.  No consistency among different estimates of
the same diagonal entry is assumed.  In the diagonal coordinates of
$A$, the form $g_{N,A}$ is the sum of the nonnegative coordinate
terms displayed above.  Each coordinate used in $B-E$ occurs in at
least one of the pair forms, while unused duplicate diagonal terms
can be discarded.  Therefore
\begin{equation}\label{eq:matrixlearningassembly77}
 g_{N,A}(B-E)^2
 \le\sum_{i<j}
       \gamma_{ij}(\widehat C_{ij}-C_{ij})^2
 \le14sA_0^2\varepsilon^2.
\end{equation}

The assembled matrix need not be legal.  Define the estimate by
minimizing its distance to the legal effect body in the known
positive quadratic norm $g_{N,A}$.  The minimum is attained and is
unique.  An exact minimizer $F$ satisfies
\[
 g_{N,A}(F-E)\le2g_{N,A}(B-E),
\]
because $E$ is a feasible competitor.  To obtain a finite procedure
with a rational legal output, it suffices instead to choose
$F\in\mathfrak E_d$ with rational entries such that
\begin{equation}\label{eq:matrixlearningprojection77}
 g_{N,A}(F-B)
 \le\inf_{G\in\mathfrak E_d}g_{N,A}(G-B)+\varepsilon.
\end{equation}
This changes the bound only to
\begin{equation}\label{eq:matrixlearningprojectionerror77}
 g_{N,A}(F-E)
 \le2\sqrt{14s}A_0\varepsilon+\varepsilon
 \le C_0d\varepsilon
\end{equation}
for an absolute $C_0$.

For completeness, \eqref{eq:matrixlearningprojection77} is an
effective finite selection, not a query to the unknown effect.
Since $W_A\succeq\tau I/2$,
$g_{N,A}(X)\le N\|X\|_{\rm HS}$.  A finite rational grid covering
$\mathfrak E_d$ in Hilbert--Schmidt norm to
$\varepsilon/(4N)$ therefore covers it in $g_{N,A}$ to
$\varepsilon/4$.  Such a legal grid is obtained by rounding after
the inward move $G\mapsto(1-2r)G+rI$: entrywise rounding with
Hilbert--Schmidt error at most $r/2$ preserves positivity of both
outcomes, and the total rounding displacement is at most $C_dr$.
Concretely, choose positive rational
$r\le\varepsilon/[4N(\sqrt d+1)]$ and a rational entry mesh at
most $r/(2d)$.  The inward displacement is at most $r\sqrt d$,
and the rounding displacement is at most $r/2$, so their sum is
at most $\varepsilon/(4N)$.
The grid and its rational entries are defined in the fixed public
input basis, independently of the random eigenbasis used for
calibration.  Enumerate this bounded grid, retain legal matrices
using exact positivity tests, and evaluate its finitely many
objective values to tolerance $\varepsilon/4$.  Selecting the
smallest computed value proves
\eqref{eq:matrixlearningprojection77} with slack.  The calibration
matrices are algebraic, and the qubit subroutine's outputs are finite
computable expressions in its record, so these evaluations require
no oracle for the unknown effect.  Only the observed matrices, the
public parameters, and finite classical calculations are used.
Arithmetic and state preparation are not included in the device-call
bound.

Finally, \eqref{eq:matrixcalibratedvariance77} gives
\[
 g_{N,E}(F-E)\le\sqrt{12}\,g_{N,A}(F-E)
                  \le\sqrt{12}C_0d\varepsilon.
\]
Choose the absolute $c_*$ in
\eqref{eq:matrixlearningpairaccuracy77} so that this is at most
$\delta/(8\sqrt2)$, and reduce $\delta_0$ if necessary so that
it is also at most $1/4$.  The endpoint-to-midpoint comparison in
Lemma~\ref{lem:matrixlocal76} and the upper bound in
Theorem~\ref{thm:matrixmetric76} then give
\[
 d_N(\mathcal M_E,\mathcal M_F)
 \le8Q_N(E,F)\le8\sqrt2\,g_{N,E}(F-E)\le\delta.
\]
This proves the risk bound.

The calibration costs at most $Cd^4N\log(d/\eta)$ calls.  There
are $s\le d^2/2$ qubit learners, each with accuracy
$c_*\delta/d$ and confidence $\eta/(2s)$.  Their combined call
count is at most $Cd^4N\delta^{-2}\log(d/\eta)$ by
Theorem~\ref{thm:commonlearn76}.  All budgets hold on every record,
all compressed entangled blocks have length at most $N$, and
postprocessing uses no additional calls.

The confidence and resource allocation can be recorded as follows;
the constants in the call bounds are absolute.
\begin{center}
\begin{tabular}{lll}
\hline
Stage & Failure allowance & Calls on every record\\
\hline
Endpoint calibration & $\eta/2$
    & $Cd^4N\log(d/\eta)$\\
Each of $s=\binom d2$ pair learners & $\eta/(2s)$
    & $CN(d/\delta)^2\log(d/\eta)$\\
Legal rational selection & $0$ & $0$\\
\hline
\end{tabular}
\end{center}
The conditional guarantees of the pair learners allow their
allowances to be added after calibration, despite the dependence of
the selected basis on its record.

For the lower bound take
\[
 E_0=\tfrac12I_d,\qquad
 E_1=(\tfrac12+\Delta)I_d,
 \qquad \Delta=256\delta/\sqrt N.
\]
Choose $\delta_0\le1/4096$.  The Bernoulli separation in
\eqref{eq:learnscalarsep76} is unchanged by the input dimension,
and the two effects are separated by more than $2\delta$ under
$d_N$.  Every use of either scalar device produces a fresh
Bernoulli bit independent of the input.  Its conditional reference
state is the same reduced state multiplied by that scalar outcome
probability.  Thus an arbitrary adaptive training record with at
most $M$ calls, padded if necessary, is a common processing of
$M$ independent Bernoulli bits and parameter-independent
resources.  Its relative entropy is at most $3M\Delta^2$ by
\eqref{eq:learnscalarkl76}.  The two disjoint success balls turn
the estimated effect into a test with both errors at most $\eta$.
The same fidelity inequality used in the proof of
Theorem~\ref{thm:commonlearn76} gives
\[
 2\eta\ge\tfrac12e^{-3M\Delta^2},\qquad
 M\ge\frac{N}{3\cdot256^2\delta^2}\log\frac1{4\eta}.
\]
For $\eta\le1/8$ this is
\eqref{eq:matrixlearninglower77}.
\end{proof}

The common learner now has the same dimension and boundary scope
as the metric and entropy theorems.  Its calibration is a sequence
of product experiments, while the compression learners provide the
enhanced angular resolution when an estimated two-dimensional
subspace is close to opposite support endpoints.  No pair of
candidate hypotheses is supplied to either part of the procedure.
